% Generated draft: requires mathematical and source review.
\documentclass[11pt]{article}
\usepackage[margin=25mm]{geometry}
\usepackage{fontspec}
\setmainfont{Latin Modern Roman}
\usepackage{amsmath,amssymb,mathtools,unicode-math}
\setmathfont{Latin Modern Math}
\usepackage{enumitem,needspace,xurl,hyperref}
\setlength{\parindent}{0pt}
\setlength{\parskip}{5pt}
\setlength{\emergencystretch}{3em}
\newcommand{\N}{\mathbb{N}}
\newcommand{\Z}{\mathbb{Z}}
\newcommand{\Q}{\mathbb{Q}}
\newcommand{\R}{\mathbb{R}}
\newcommand{\C}{\mathbb{C}}
\newcommand{\F}{\mathbb{F}}
\newcommand{\A}{\mathbb{A}}
\newcommand{\PP}{\mathbb{P}}
\newcommand{\E}{\mathbb{E}}
\newcommand{\eps}{\varepsilon}
\newcommand{\dd}{\,\mathrm{d}}
\begin{document}
\section*{UnsolvedMath ID: 30002708}
\textbf{Source problem code:} catalogue entry
\subsection{Definitions and mathematical statement}
The supplied record asks for a motivic analogue of Kontsevich--Zagier equivalence, but does not specify enough data to determine a unique mathematical conjecture. The following is a statement schema, with its unresolved parameters made explicit.

Let \(\mathcal P\) be a specified class of admissible integral presentations \(P=(X,f)\) in a Cluckers--Loeser motivic integration theory, with \(f\) integrable on \(X\). Fix a common target \(V\) for the integral values and write
\[
I\colon\mathcal P\longrightarrow V,\qquad I(X,f)=\int_X f.
\]
Let \(\mathcal R\subseteq\mathcal P\times\mathcal P\) consist of the admissible elementary relations expressing the intended motivic analogues of additivity, semialgebraic change of variables, and Stokes' formula. Under the interpretation that ``related by those rules'' means the equivalence relation generated by these elementary relations, define \(P\sim_{\mathcal R}Q\) to mean that there exist \(n\in\Z_{\geq 0}\) and \(P_0,\ldots,P_n\in\mathcal P\) with \(P_0=P\), \(P_n=Q\), and, for every \(0\leq j<n\),
\[
(P_j,P_{j+1})\in\mathcal R
\quad\text{or}\quad
(P_{j+1},P_j)\in\mathcal R.
\]
The requested implication is
\[
\forall\,(X_1,f_1),(X_2,f_2)\in\mathcal P,\qquad
\int_{X_1}f_1=\int_{X_2}f_2
\ \Longrightarrow\
(X_1,f_1)\sim_{\mathcal R}(X_2,f_2).
\]
Here the finite, reversible interpretation of rule application is an explicit formalization convention requiring confirmation. The supplied original statement leaves the classes of sets, functions, and rules unspecified. The clean statement identifies Cluckers--Loeser motivic integrals and names three kinds of rules, but supplies neither their precise versions nor the underlying theory, admissibility hypotheses, boundary conditions, or meaning of equality. Consequently, \(\mathcal P\), \(V\), and \(\mathcal R\) remain unresolved parameters; the displayed implication is not a fully specified conjecture.

The supplied status is \emph{partially solved}. Its accompanying assessment concerns the broad real period problem and a positive \(p\)-adic semialgebraic analogue. The record does not explain whether or how that result settles a case of the stated motivic implication.
\subsection{Short English statement}
The question asks whether every equality between two admissible Cluckers--Loeser motivic integrals can be explained by the specified motivic versions of additivity, semialgebraic change of variables, and Stokes' formula. The admissible integrals, equality of values, and exact rules remain unspecified. The supplied partial-solution assessment concerns related real and \(p\)-adic problems and does not establish a result for this motivic formulation.
\subsection{Sources}
\begin{itemize}
\item Local UnsolvedMath registry, supplied complete record for permanent ID \(30002708\), title ``Motivic Kontsevich--Zagier Equivalence,'' exact source problem code \texttt{catalogue entry}, in the Oberwolfach Reports Open Problems collection. The supplied proposed year is 2014; no individual proposer is named.
\item Supplied report citation: \emph{Valuation Theory and Its Applications} (2015), Oberwolfach Reports, DOI 10.4171/owr/2014/49, pp.~2757--2823. \url{https://doi.org/10.4171/owr/2014/49}. The background describes Oberwolfach Reports as records of mathematical workshops and their open-problem discussions.
\item Supplied literature source: Cluckers--Halupczok, ``\(p\)-adic Kontsevich--Zagier rules.'' \url{https://arxiv.org/abs/2003.01500}.
\item Supplied literature source: ``Oberwolfach Report 49/2014.'' \url{https://doi.org/10.4171/owr/2014/49}.
\item Literature checked at: \texttt{2026-08-21}. This is supplied dated curation, not an independently verified assessment. The supplied assessment states that the broad real period problem remains open, while Cluckers and Halupczok prove a positive \(p\)-adic semialgebraic analogue using basic integral transformation rules. The background explicitly calls this dated literature triage derived from the supplied OWR curation and requests independent verification of time-sensitive status and cited claims.
\item The supplied statement status is \texttt{corrected\_verified}. The accompanying registry assertion says that the statement was checked against the official Oberwolfach report and preserves its hypotheses, notation, and open target. That assertion is source metadata and has not been independently verified in this draft.
\end{itemize}
\end{document}
